% !TEX root = ../risk_vs.tex
\section{Heterogeneous Research and Outperformance}\label{sec:hetero}

Peer-reviewed publications differ in many dimensions, including theoretical foundations, discipline of origin, and outlet quality. This section examines these differences and whether they matter for predictive outperformance compared to data mining.


% == data summary == #
\subsection{Research Categorization Methods\label{sec:hetero-method}}

We categorize research at the predictor-paper level, along four dimensions:
\begin{enumerate}
   \item Theoretical foundation: is the theoretical explanation for the predictor risk or mispricing? Or is there no clear theoretical explanation (agnostic)?
   \item Equilibrium modeling: is the theoretical explanation supported by a stylized model, dynamic model, or quantitative model?
   \item Academic Discipline: is the paper published in a finance, accounting, or other discipline's journal?
   \item Quality: is the research published in a top-3 finance or top-3 accounting journal?
\end{enumerate}
These dimensions are key characteristics of asset pricing research. Conferences and societies are organized around the first three. Journal quality is important for tenure. 

We apply these classifications based on manual reading of the original papers. All classifications can be found at \url{https://github.com/chenandrewy/flex-mining/blob/main/DataInput/SignalsTheoryChecked.csv}, which also contains quotes justifying classifications for the first two dimensions (theory and equilibrium modeling).

We categorize predictor-papers, rather than papers, because similar predictors appear across multiple papers, and the peer review process sometimes arrives at different theoretical foundations. Similar predictors appear in the CZ dataset, on occasion, due to their goal of comprehensive coverage and the fact that it's difficult to judge if two predictors are indeed duplicates.

For example, predictors related to profitability appear in both \citet{fama2006profitability} and \citet{balakrishnan2010post}. Fama and French use annual earnings and scale by book equity, while Balakrishnan et al. use quarterly earnings and scale by total assets.\footnote{% begin footnote
While one might argue that these predictors are duplicates, Balakrishnan et al say their predictor is ``incremental to, and more pronounced than previously documented earnings-related anomalies.''
} % end footnote 
Fama and French provide an agnostic explanation: ``We take no stance on whether the patterns in average returns observed here are rational or irrational.'' In contrast, Balakrishnan et al. argue for mispricing: ``We document a market failure to fully respond to loss/profit quarterly announcements.'' Thus, the broader concept of profitability appears as both agnostic and mispricing in our analysis, and is evaluated at the predictor-paper level. 

In rare instances, predictors from the same paper may have distinct theoretical foundations. \citet{ang2006cross} provide a risk-based explanation for the predictive power of the VIX beta (``innovations in aggregate volatility carry a statistically significant negative price of risk''), but are agnostic about why idiosyncratic volatility predicts returns (``our results on idiosyncratic volatility represent a substantive puzzle''). Thus, the Ang et al. paper appears as both risk-based and agnostic, and one needs to delve into the predictor-paper level to have a clean description of the theoretical ideas.

As these quotes illustrate, categorizing predictor-papers into risk, mispricing, and agnostic is typically straightforward. However, a handful of predictor-papers focus on liquidity mechanisms.  We categorize these predictor-papers as mispricing if the argument focuses on stock-specific measures of liquidity (\citealt{amihud2002illiquidity})  and risk if the argument focuses on a market-wide component (\citealt{pastor2003liquidity}). This method gives the risk category the best chance at finding post-sample returns, since idiosyncratic liquidity has improved over time (\citealt{chen2022zeroing}). Nevertheless, this issue affects only seven predictor-papers, and has little impact on our main results.


\subsection{The Distribution of Research Categories}\label{sec:hetero-distribution}

Table \ref{tab:hetero-count} illustrates the distribution of predictor-papers across the four categories. Among the 173 predictor-papers we examine, 20\% (34/173) are attributed to risk by the peer review process, 61\% (105/173) are attributed to mispricing, and 20\% (34/173) have uncertain or agnostic origins. Top finance journals show a similar pattern, with 22\% (23/105) of predictors attributed to risk. Interestingly, accounting journals rarely attribute predictability to risk. A detailed breakdown by journal is in the Internet Appendix (Table \ref{tab:journal_theory}).

% table: research approach summary 
\begin{table}[!h]\caption{Research Categories in Cross-Sectional Predictability}
\label{tab:hetero-count}

\begin{singlespace}
\noindent  We categorize predictor-papers by reading the original papers. Justification for each ``Theoretical Explanation'' is at \url{https://github.com/chenandrewy/flex-mining/blob/main/DataInput/SignalsTheoryChecked.csv}. ``JF, JFE, RFS'' includes only predictors published in the Journal of Finance, Journal of Financial Economics, or Review of Financial Studies. ``AR, JAE, JAR'' includes only predictors published in the Accounting Review, the Journal of Accounting and Economics or the Journal of Accounting Research. Peer review attributes little cross-sectional predictability to risk. Few predictor-papers are explained with an equilibrium model.
\end{singlespace}
\begin{centering}
\vspace{-1ex}
\par\end{centering}
\centering{}\setlength{\tabcolsep}{1ex} \footnotesize
\begin{center}
   \begin{tabular}{lrrrrr}
      \toprule
      \multicolumn{1}{l}{\textbf{Category}} 
        & \multicolumn{5}{c}{\textbf{Journal Category}} \\
      \cmidrule(lr){2-6}
      & Any & JF, JFE, RFS & AR, JAE, JAR & Finance & Accounting \\
      \midrule
      \multicolumn{6}{l}{\textbf{Theoretical Explanation}} \\
      % latex table generated in R 4.4.1 by xtable 1.8-4 package
% Mon Dec  8 10:27:44 2025
 Risk &  34 &  23 &   0 &  28 &   1 \\ 
  Mispricing & 105 &  59 &  24 &  68 &  32 \\ 
  Agnostic &  34 &  23 &   5 &  23 &   5 \\ 
  
      \multicolumn{6}{l}{\textbf{Equilibrium Modeling}} \\
      % latex table generated in R 4.4.1 by xtable 1.8-4 package
% Mon Dec  8 10:30:02 2025
 No Model & 148 &  87 &  29 &  98 &  38 \\ 
  Stylized &  14 &  10 &   0 &  12 &   0 \\ 
  Dynamic &   5 &   4 &   0 &   5 &   0 \\ 
  Quantitative &   6 &   4 &   0 &   4 &   0 \\ 
  
      \midrule
      % latex table generated in R 4.4.1 by xtable 1.8-4 package
% Mon Dec  8 10:31:09 2025
 \textbf{Total} & 173 & 105 &  29 & 119 &  38 \\ 
  
      \bottomrule
   \end{tabular}
\end{center} 
\end{table}

This distribution suggests a consensus about the origins of cross-sectional predictability: peer review attributes little of this predictability to risk. We emphasize that this attribution is chosen not only by the authors, but also by the editors and referees. This consensus is consistent with factor model measures of risk (Appendix \ref{sec:intapp-risk-factor}). It contrasts with recent reviews of empirical asset pricing, which are typically agnostic about the origins of predictability (e.g. \citealt{bali2016empirical,zaffaroni2022asset}). 

The distribution also shows that relatively few predictor-papers are justified with an equilibrium model. Only 14.5\% of the predictors-papers have a mathematical theory of any sort. And among these models, most are stylized. This result may be related to the high prevalence of mispricing explanations, which can be thought of as disequilibrium phenomena. Indeed, of the 25 papers with a mathematical theory, only five are theories of mispricing, and most of these center around trading frictions. Alternatively, this result may be due to the technical challenges around solving equilibrium models with a cross-section of stocks under uncertainty. 

The papers that lack a mathematical theory vary in the assertiveness of their theoretical explanations. Some papers begin with a verbal theory of investor behavior (sometimes with citations to mathematical theories), and then present empirical evidence consistent with the theory (e.g. \citealt{sloan1996stock}; \citealt{asquith2005short}; \citealt{bali2011maxing}; \citealt{eberhart2004examination}). Others take a more agnostic stance, discussing several possible theories, and then use empirical evidence to argue for one of the theories (\citealt{spiess1999long}; \citealt{titman2004capital}; \citealt{chan1996momentum}). 80\% of predictor-papers end up with a stance on the theoretical origin of predictability, as seen in Table \ref{tab:hetero-bytheory}. 

\subsection{Post-Sample Outperformance by Theoretical Support}\label{sec:hetero-mispricing}

% Describe the table's methods
Table \ref{tab:hetero-bytheory} repeats the post-sample analysis from Figures \ref{fig:intro} and \ref{fig:pub-vs-dm-2}, applied to subsets of predictor-papers based on their theoretical foundation and equilibrium modeling.

% table: outperformance by research approach
\begin{table}[!h]
   \caption{Post-Sample Outperformance by Theoretical Support}
   \label{tab:hetero-bytheory}
   
   \begin{singlespace}
      \noindent Strategies are normalized to have 100 bps per month performance in-sample. `Post-Sample' is the performance of predictor-papers in bps per month. `Versus Data Mining' is `Post-Sample' minus the performance of data-mined benchmarks. `Theoretical Foundation' and `Equilibrium Modeling' are determined by reading the papers (see Table \ref{tab:hetero-count}). CAPM- and FF3+momentum-alphas use regressions specific to the sample periods (in-sample or post-sample). Standard errors (parentheses) are clustered by calendar month and predictor. Research that is agnostic on the origin of predictability shows signs of outperformance relative to data mining. Research that takes a stand on the theory does not.
   \end{singlespace}
   \begin{centering}
   \vspace{0ex}
   \par\end{centering}
   \centering{}\setlength{\tabcolsep}{1.2ex} \small
   \begin{center}
   \begin{tabular}{lcccccc}
      \toprule
      & \multicolumn{2}{c}{Long-Short Return} & \multicolumn{2}{c}{CAPM Alpha} & \multicolumn{2}{c}{FF3 + Mom Alpha} \\
      \cmidrule(lr){2-3}\cmidrule(lr){4-5}\cmidrule(lr){6-7}
      & Post- & \multicolumn{1}{c}{Versus} & Post- & \multicolumn{1}{c}{Versus} & Post- & \multicolumn{1}{c}{Versus} \\
      & Sample & \multicolumn{1}{c}{Data Mining} & Sample & \multicolumn{1}{c}{Data Mining} & Sample & \multicolumn{1}{c}{Data Mining} \\
      \midrule
      % FF4 figures copied from exhibits/ff4/Table_RiskAdjusted_ff4_t2.tex
      \multicolumn{7}{l}{\textbf{Theoretical Foundation}} \\
      Agnostic          & 65 & 9 & 79 & 23 & 110 & 31 \\
                        & (8) & (8) & (8) & (8) & (8) & (9) \\
      Mispricing        & 55 & 4 & 60 & 6 & 59 & -17 \\
                        & (4) & (4) & (4) & (4) & (4) & (5) \\
      Risk              & 43 & 5 & 38 & -4 & 49 & -21 \\
                        & (8) & (8) & (8) & (8) & (6) & (9) \\                        
      \addlinespace
      % FF4 figures copied from exhibits/ff4/Table_RiskAdjusted_ff4_t2.tex
      \multicolumn{7}{l}{\textbf{Equilibrium Modeling}} \\
      No Model          & 56 & 5 & 62 & 9 & 71 & -4 \\
                        & (3) & (3) & (3) & (3) & (3) & (4) \\
      Stylized          & 63 & 15 & 49 & -5 & 51 & -42 \\
                        & (12) & (13) & (13) & (14) & (7) & (11) \\
      Dynamic or        & 34 & -2 & 50 & 4 & 39 & -54 \\
      Quantitative      & (14) & (14) & (14) & (19) & (14) & (28) \\
      \addlinespace
      \textbf{Overall}  & 56 & 5 & 60 & 8 & 68 & -8 \\
                        & (3) & (3) & (3) & (3) & (3) & (4) \\
      \bottomrule
   \end{tabular}
   \end{center} 
\end{table}

% walk through results. begin with the simple pattern that agnostic performs best
The top panel shows that research with agnostic theoretical foundations has the strongest post-sample performance. As seen in the `Post-Sample' columns, agnostic research produces 65 to 110 bps per month post-sample, depending on the performance metric. In comparison, mispricing foundations produce 55 to 60 bps per month post-sample, and risk foundations produce 38 to 49 bps.

% results: main question - does anything outperform data-mining?
The ``Versus Data Mining'' columns report the difference between the post-sample performance of the predictor-papers and their data-mined benchmarks, and thereby address our main question: does peer-reviewed research with a particular theoretical foundation outperform data mining?

Research with theoretical foundations in mispricing or risk show little to no outperformance. Predictor-papers with mispricing foundations outperform their data-mined benchmarks by only 4 bps per month in terms of long-short returns and 6 bps in terms of CAPM alphas. In terms of FF3 + momentum alphas, mispricing foundations \emph{under}perform data mining, by 17 bps per month. The performance of predictor-papers with risk foundations is even worse. By most metrics, risk-based predictor-papers underperform data mining.

% main result: agnostic methods outperform
Agnostic theoretical foundations show some signs of outperformance, but the magnitudes are modest. Agnostic predictor-papers outperform data mining by 9 to 31 bps per month, depending on the performance metric. Since performance is normalized to 100 bps in-sample, this means agnostic predictors retain up to an additional 31 percentage points of their original-sample performance, compared to data mining. While 31 percentage points may appear notable, it is the largest estimate out of many, and should be shrunk toward the average of about zero to account for multiple comparisons (e.g. \citealt{chen2020publication}). This finding, that agnostic research outperforms somewhat while risk-based and mispricing-based research do not, is robust to alternative construction methods of the data-mined benchmarks, including directly controlling for t-stats and mean returns of published strategies (Appendix \ref{sec:app-robust}). 

% results: equilibrium modeling
Grouping research by equilibrium modeling leads to a similar pattern, as seen in the bottom panel of Table \ref{tab:hetero-bytheory}. By most metrics, predictor-papers with no equilibrium model have stronger post-sample performance than predictor-papers with equilibrium modeling. The performance of model-free research is not as strong as research with agnostic theoretical foundations, however. Predictor-papers with no model produce 56 to 71 bps per month post-sample, compared to 65 to 110 bps per month for agnostic predictor-papers. Indeed, when compared to data mining, model-free research shows minimal outperformance. 

Notably, Table \ref{tab:hetero-bytheory} implies that the support of a stylized or dynamic equilibrium model leads to little if any outperformance relative to data mining. Additional robustness checks are in the Internet Appendix (\ref{sec:intapp-riskalt} and \ref{sec:intapp-pub-additional}).


\subsection{Post-Sample Outperformance by Discipline and Journal Ranking}\label{sec:hetero-discipline}

Table \ref{tab:hetero-byjournal} examines how discipline and journal ranking affect outperformance. The top panel shows that predictor-papers in finance journals have stronger post-sample performance compared to accounting journals. Finance predictor-papers produce 59 to 76 bps per month post-sample, compared to 43 to 56 bps per month for accounting journals. 

\begin{table}[!h]
   \caption{Post-Sample Outperformance by Discipline and Journal Ranking}
   \label{tab:hetero-byjournal}
   
   \begin{singlespace}
   \noindent Strategies are normalized to have 100 bps per month performance in-sample. `Post-Sample' is the performance of predictor-papers in bps per month. `Versus Data Mining' is `Post-Sample' minus the performance of data-mined benchmarks. `Discipline' categorizes the journal of publication. `JF,' `JFE,' and `RFS' includes papers in the Journal of Finance, Journal of Financial Economics, and Review of Financial Studies. `AR,' `JAR,' and `JAE' includes papers in the Accounting Review, Journal of Accounting Research, and Journal of Accounting and Economics. CAPM and FF3 + momentum alphas use regressions specific to the sample periods (in-sample or post-sample). Standard errors (parentheses) are clustered by calendar month and predictor. Finance, and particularly top-ranked finance journals show some signs of outperformance relative to data mining, but the improvement is modest. 
   \end{singlespace}
   \begin{centering}
   \vspace{0ex}
   \par\end{centering}
   \centering{}\setlength{\tabcolsep}{1.2ex} \small
   \begin{center}
   \begin{tabular}{lcccccc}
      \toprule
      & \multicolumn{2}{c}{Long-Short Return} & \multicolumn{2}{c}{CAPM Alpha} & \multicolumn{2}{c}{FF3 + Mom Alpha} \\
      \cmidrule(lr){2-3}\cmidrule(lr){4-5}\cmidrule(lr){6-7}
      & Post- & \multicolumn{1}{c}{Versus} & Post- & \multicolumn{1}{c}{Versus} & Post- & \multicolumn{1}{c}{Versus} \\
      & Sample & \multicolumn{1}{c}{Data Mining} & Sample & \multicolumn{1}{c}{Data Mining} & Sample & \multicolumn{1}{c}{Data Mining} \\
      \midrule
      % FF4 figures copied from exhibits/ff4/Table_RiskAdjusted_DisciplineJournal_ff4_t2.tex
      \multicolumn{7}{l}{\textbf{Discipline}} \\
      Finance      & 59 & 8 & 65 & 12 & 76 & -2 \\
                   & (4) & (4) & (4) & (4) & (4) & (5) \\
      Accounting   & 43 & -6 & 45 & -8 & 43 & -27 \\
                   & (6) & (7) & (6) & (6) & (5) & (7) \\
      \addlinespace
      \multicolumn{7}{l}{\textbf{Journal Rank}} \\
      JF, JFE, RFS & 60 & 8 & 68 & 16 & 81 & 3 \\
                   & (4) & (4) & (4) & (5) & (4) & (5) \\
      AR, JAR, JAE & 43 & -6 & 45 & -8 & 43 & -27 \\
                   & (6) & (7) & (6) & (6) & (5) & (7) \\
      Other        & 53 & 8 & 55 & 2 & 61 & -20 \\
                   & (6) & (6) & (6) & (7) & (7) & (9) \\
      \bottomrule
   \end{tabular}
   \end{center} 
\end{table}

This performance is less impressive when compared to data mining, however. Predictor-papers in finance journals outperform data mining by between -2 and +12 bps per month.  Predictor-papers in accounting journals underperform data mining, regardless of the performance metric.

The top 3 finance journals (Journal of Finance, Journal of Financial Economics, Review of Financial Studies) perform a bit better than finance journals overall. Nevertheless, the outperformance relative to data mining is small. Predictor-papers in these top journals outperform data mining by 3 to 16 bps per month, depending on the performance metric. Given the normalization, this means top finance journals retain an additional 3 to 16 percentage points of their original-sample performance, relative to simple data mining. 